% !PRINTED 30
\weberpage{30}

\begingroup
\setlength{\abovedisplayskip}{5pt plus 2pt minus 1pt}
\setlength{\belowdisplayskip}{5pt plus 2pt minus 1pt}
\setlength{\abovedisplayshortskip}{3pt plus 1pt}
\setlength{\belowdisplayshortskip}{4pt plus 1pt}
\linespread{0.94}\selectfont

\begin{center}
{\large \S.\ 4.}

\vspace{0.5em}
{\large\gesperrt{Division by a Linear Function.}}
\end{center}

We shall apply the general rule for division to one special case. Let
the dividend be
\begin{equation}
 f(x)=a_0x^n+a_1x^{n-1}+a_2x^{n-2}+\cdots+a_n
\tag{1}
\end{equation}
arbitrary, while the divisor is of the first degree, or, as one also
says, linear. We shall also suppose the coefficient of the first power
of $x$ in the divisor to be $=1$, and therefore take the divisor in the
simple form $(x-\alpha)$, where $\alpha$ remains arbitrary. In this
case the quotient $Q$ is of degree $(n-1)$, and the remainder $C$ is of
degree $0$, that is, independent of $x$. Thus set
\begin{equation}
 f(x)=(x-\alpha)Q+C,
\tag{2}
\end{equation}
so that $C$ no longer contains $x$; and if we set
\begin{equation}
 Q=q_0x^{n-1}+q_1x^{n-2}+\cdots+q_{n-2}x+q_{n-1}
\tag{3}
\end{equation}
then there follows from (2):
\begin{equation}
\begin{aligned}
 f(x)={}&q_0x^n+q_1x^{n-1}+\cdots+q_{n-2}x^2+q_{n-1}x+C\\
 &-\alpha q_0x^{n-1}-\cdots-\alpha q_{n-3}x^2
   -\alpha q_{n-2}x-\alpha q_{n-1},
\end{aligned}
\tag{4}
\end{equation}
and by comparison with (1):
\begin{equation}
\begin{aligned}
 q_0&=a_0,\\
 q_1-\alpha q_0&=a_1,\\
 q_2-\alpha q_1&=a_2,\\
 &\cdots,\\
 q_{n-1}-\alpha q_{n-2}&=a_{n-1},\\
 C-\alpha q_{n-1}&=a_n.
\end{aligned}
\tag{5}
\end{equation}
From this one obtains
\begin{equation}
\begin{aligned}
 q_0&=a_0,\\
 q_1&=a_0\alpha+a_1,\\
 q_2&=a_0\alpha^2+a_1\alpha+a_2,\\
 &\cdots,\\
 q_{n-1}&=a_0\alpha^{n-1}+a_1\alpha^{n-2}+\cdots+a_{n-1},\\
 C&=a_0\alpha^n+a_1\alpha^{n-1}+\cdots+a_{n-1}\alpha+a_n=f(\alpha).
\end{aligned}
\tag{6}
\end{equation}
$C$ arises from $f(x)$ by setting $x=\alpha$, and therefore can also
be denoted by $f(\alpha)$.
\endgroup
